\section{Use available primary depth to share the order-twelve square}

The frozen preceding source has SHA256
\texttt{78b54bcde24ff8a91091e102a4734c49878699f071d9230e5ce7a70ffd2114d0}.
The isolated replacement in this section has SHA256
\texttt{5eae7a6cb6fd22e2da8a457354628355ece7e8930b60200d088cd0cfca1d28cb}.
It saves one rational coordinate in each of the three detectors by
strengthening one test which the existing completeness witnesses already
satisfy to higher order. No global arithmetic input is proved here; their
qualifications and the preceding complete local arguments are retained.

\subsection{The actual test whose order changes}

Write $U_A=[u_A]P_*$ for the primary point attached to the input label,
$g_A=g(U_A)$ for its complete reciprocal-chart coordinate, and
\[
 D_{\times}=L_1-\overline{\ell(V_1)/\ell(V_0)}
\]
for the existing input cross test. Thus $g_A^2=f(U_A)$, and the source
aliases are \texttt{test1} and \texttt{test10}, respectively. The
preceding detector required orders one and twelve on this pair. Require
order twelve on \emph{both} tests instead. The complete order list, in
the source order \texttt{test0},\ldots,\texttt{test15}, is now
\[
 (45,12,1,1,1,1,1,1,1,1,12,45,45,45,45,45).
\]
There are eight order-one tests, two order-twelve tests, and six
order-forty-five tests. All underlying field functions, their conjugations,
and all other orders are unchanged. In particular no order-one alternate
point test is strengthened: the alternate points still need only formal
depth one.

The retained one-place lemma applies to arbitrary positive orders.
With $t=A\lambda$, $\delta=98t-1$, and
\[
 R=\frac{F(t)}{A^5\delta^6},\qquad
 q_m=\frac{e^{2m-1}}{\delta},\qquad
 e=\frac{(A\delta)^5}{F(t)},
\]
the semantic test is still $\delta D$, and the implemented factor is
$1+q_mD^2$. At a consumed contact $\delta$ is a unit, so the common-place
conclusion for the new order list supplies $v_{\mathfrak p}(g_A)\ge12$.
This implies the former order-one conclusion and leaves every other
required conclusion unchanged. The subsequent primary/alternate
comparison, the actual $G_2/G_{12}$ error bounds, the identity-input case,
and the three height consumers therefore apply with their old premises.
This is a strengthened local predicate used in the conditional soundness
argument; it is not a bijection between every pair of old and new detector
witness tuples.

\subsection{Why completeness already gives this stronger primary test}

In the retained completeness construction, the integer input $X$ first
fixes the integer label
\[
 n=\beta(X)=11016X^2+11016X+5155.
\]
The primary points are chosen as $H=[h]P_*$ and $U_A=[hn]P_*=[n]H$.
At every selected contact place used for the primary side, the existing
progression $\Pi^{44}\mid h$, together with the prescribed annihilating
reduction order, makes $H$ formal of depth at least forty-five. This is
exactly the existing depth used for the high-order \texttt{test0}; the
present change imposes no new congruence on $h$.

For completeness of this deduction, let $E_s$ be the subgroup of a good
integral formal group whose parameter $z$ has valuation at least the
positive integer $s$. The integral formal addition and inverse series
show that $E_s$ is a group. Hence $[n]E_s\subseteq E_s$ for every integer
$n$, including negative $n$ and $n=0$. Equivalently the integral series
$z([n]Q)=n z(Q)+O(z(Q)^2)$ cannot lower its positive valuation. Therefore
\[
 v_{\mathfrak p}(z(U_A))\ge
 v_{\mathfrak p}(z(H))\ge45.
\]
On a formal point in the reciprocal chart,
$z=-g/l$ and $l^2=1-Dg^4$; the relevant $l$ is a unit. Consequently
\[
 v_{\mathfrak p}(g_A)=v_{\mathfrak p}(z(U_A))\ge45\ge12.
\]
Both choices of the square-root chart have the same valuation. The
identity case has $g_A=0$ and satisfies every requested positive order;
no nonidentity guard for $U_A$ has been inserted. The consumed primes
remain outside $\{2,3,5,7,13,17\}$, so the good integral and unramified
formal-group hypotheses are exactly those already used in the retained
primary-depth argument.

The existing CRT conditions and their coprimalities are unchanged.
In particular the alternate condition involving $h-1$ is not strengthened,
and one does not demand that a selected prime divide both $h$ and $h-1$.
The scalar preparation and continuity neighborhoods, root points, square
scales and finite valuation/real bounds are still fixed before the
asymptotic contact-prime selection. This argument uses the actual
completeness choice $U_A=[n]H$; it does not assume that equality for an
arbitrary putative solution in the soundness direction.

\subsection{The shared uniform block and its exact cost}

Put $U=g_A^2$ and $V=D_{\times}^2$, both elements of $K=\mathbb Q(\sqrt2)$.
The preceding product block was
\[
 (1+q_1U)(1+q_{12}V)
 =1+q_1U+q_{12}V+(q_1q_{12})UV.
\]
The replacement is the uniform-order block
\[
 (1+q_{12}U)(1+q_{12}V)
 =1+q_{12}(U+V)+q_{12}^2 UV.
\]
The shared field coefficient $UV$ is exactly the same as before.
The stored square $q_{12}^2$ is already the coordinate
\texttt{detector*.mixedHigh}, required by the unchanged high-weight
relation
\[
 q_1^3q_{45}=(q_{12}^2)^2.
\]
Thus each detector can delete its separate coordinate for $q_1q_{12}$.
The two rational output coordinates of the block remain; every defining
row is quadratic in the stored coordinates. The two four-test order-one
blocks and two three-test high-order blocks use exactly their preceding
lists, because the promoted test was precisely the one previously
removed from the order-one list for the mixed pair. Their shared
coefficients, product evaluations and weight powers are unchanged.
The seven-node higher-weight chain also remains unchanged.

On the positive source domain $F,q_1,q_{12},q_{45},E$ are positive.
Both squared tests are nonnegative at either real embedding. The new
block and each other factor are therefore at least one. The old
conclusions $W_0\ge1$, $N(W)>0$, nonzero reciprocal trace, and positive
pole input remain valid. No quotient denominator or branch restriction
has been added. In the construction of the weight chain, all denominators
remain the previously checked nonzero $E$ and $q_1^3$.

\subsection{The changed factor in the pole-completeness estimates}

At a selected contact prime, $v_p(R)=1$ and $\delta$ is a unit. At the
successful place, $v(g_A)\ge45$ gives
\[
 v(q_{12}g_A^2)=1-24+2v(g_A)\ge67>0.
\]
The new factor is thus one modulo the prime. At either place, a
nonzero nonconstant term has odd valuation. Each inverse factor has
residue zero or one, and the reciprocal trace has the same unit residue
one or two as in the preceding proof. At noncontact primes the even
valuation of the unchanged $R$ supplies the same square-denominator
argument for the pole input.

At real infinity every fixed-order weight still satisfies
$q_m=O(\lambda^{-1})$. The underlying $g_A$ is one of the already
uniformly bounded tests in the real progression argument, so replacing
$q_1g_A^2$ by $q_{12}g_A^2$ changes only its fixed bound. The full product
still satisfies $W=1+O(\lambda^{-1})$ and the same pole input tends to
zero through positive values. In particular the required bound $B<17$
is eventually satisfied.

At $p=2,3,17$, retain
$\nu=v_p(v)$, $\eta=v_p(A)$ and $\alpha=v_p(98)$. The preceding exact
weight estimate is
\[
 v_p(q_m)=\nu+(10m-6)(\eta+\alpha).
\]
For the promoted test and its fixed lower bound $v(g_A)\ge-c_p$ this
becomes
\[
 v_p(q_{12}g_A^2)\ge
 \nu+114(\eta+\alpha)-2c_p\longrightarrow+\infty.
\]
Thus the same freedom to increase the denominator divisibility supplies
the preparation-zero pole conditions. The scale, bound $c_p$, and order
twelve are fixed before this choice. The argument does not require
$\eta+\alpha$ to be nonnegative. Together with the unchanged local
squareclasses and pole-input formula, this verifies the new
completeness application.

\subsection{Measured construction and verification scope}

With all other seventh-round parameters retained, including
$S_0=3^{2790}$, full compilation gives
\[
 400\text{ coordinates}/368\text{ residuals}
 \quad\longrightarrow\quad
 \boxed{389\text{ coordinates}/357\text{ residuals}}.
\]
The source has 339 gates. The independent gate checker verifies all
339 identities and 330 compiler operations. The compiled polynomial
has degree four, 26,295 nonzero monomials, primitive content one, and
all 389 coordinates active. These counts include the distinguished
input, which remains unchanged. The three-coordinate saving and
six-monomial decrease are relative to the frozen 392-coordinate,
360-residual construction; no final-elimination saving is counted twice.

The separately executed source checker
\texttt{check\_detector\_round8.py} imports neither producer nor compiler.
It reconstructs all sixteen semantic tests and the new order allocation,
the unchanged logarithms, weight chain, shared coefficients, uniform
pair, remaining products, norm/pole consumers and lifted degree. It
passes 384 source assertions, seven Laurent identities, three
normalization identities and two expanded component identities for the
uniform pair. These exact finite checks bind the implementation;
the formal-group and completeness arguments above supply the uniform
mathematical deduction. All inherited global arithmetic remains
conditional as in the preceding package.
